% ===== §16 =====
\section{The Poem, a Language That Circles the Gap}
\label{p26:sec:poem}

\noindent
This section opens Part Four, which draws out the generativity of the limit in three registers of relational life. Its aim is to establish that the poem is the language reconciled to its own limit and made generative by it: the one use of language that does not aim across the gap to close it but organizes itself around the gap and keeps it open, and that is for this reason the language of the drive. The section states the contrast between ordinary language and the poem, gives the account with its warrant, reviews the readings of poetic language that approach the same point, and draws the claim.

\subsection*{Ordinary language aims across the gap; the poem circles it}

Ordinary language aims at the gap in order to close it. It reaches toward the thing it would say and counts itself successful when it has said it, when the signifier has covered its object and the remainder has been, so far as possible, brought in. Its ideal is transparency, the vanishing of the medium before the message, a saying that leaves nothing unsaid. The poem does not work this way, and its difference is not a matter of ornament but of direction. The poem is organized around the gap rather than aimed across it. Its figures, its silences, its refusals of paraphrase, its way of meaning more and other than it states, all of these keep the poem turning around something it does not deliver, holding open the very remainder that ordinary speech would try to close.

\subsection*{The poem as the language of the drive}

The poem is the language of the drive, in the exact sense the paper has given the drive. The drive circles the object-cause and is satisfied in the circling, not in the seizing; the poem circles its unsayable center and is fulfilled in the circling, not in a paraphrase that would seize it. This is why a poem cannot be replaced by a statement of its content. The statement aims across the gap and closes what it can; the poem circles the gap and keeps it open, and the keeping-open is what the poem is for. To paraphrase a poem is to attempt on it exactly the closure that would, by the argument of Part Two, extinguish what made it live, and the felt deadness of a paraphrased poem is the small demonstration of that argument in the field of language itself. What the poem gives is not the closing of the gap but the experience of circling it.

\subsection*{Readings that approach the same point}

Two lines of reflection on poetic language approach this account and help to fix it. Benjamin, writing on translation, located the life of a poem in something that is not its communicable content, in a way of meaning that survives no simple transfer of information and that translation must reach toward without ever delivering as a message \parencite{benjamin1968task}. What resists translation in a poem, on this reading, is precisely what the poem circles and does not state, the unsayable center that no transfer of content carries across. Derrida, pressing the general point that no signifier delivers a final presence, that meaning is always deferred along the play of differences, showed that writing does not close upon a present meaning but keeps it in deferral \parencite{derrida1978writing}. The poem, on the present account, does not suffer this deferral as a defect but takes it up as its art, making of the non-arrival of meaning the very movement it offers. Where ordinary language treats deferral as an obstacle to be minimized, the poem treats it as the medium of what it does.

\subsection*{The core given circulation without capture}

This is why the poem writes what cannot otherwise be written. The unsymbolizable core, which no proposition reaches, is not reached by the poem either, but the poem does what the proposition cannot: it gives the core a form of circulation without giving it a form of capture. It lets the remainder be approached, turned around, dwelt near, without being seized and so extinguished. The poem is the drive's own language because it does with words what the drive does with its object, turning around a void it does not fill and generating, in the turning, something that a filling would have destroyed.

\begin{claim}[The poem circles what it cannot close]
The poem is the use of language organized around the gap rather than aimed across it. It circles an unsayable center and is fulfilled in the circling, not in a paraphrase that would seize it, and it therefore gives the unsymbolizable core a form of circulation without capture. The poem is the language of the drive, and the impossibility of paraphrasing it without deadening it is the demonstration, in language itself, that closure extinguishes what the open gap keeps alive.
\end{claim}

\begin{casebox}{the poem that says what the letter could not}
A person tries to write plainly what another means to them and finds the plain statement always beside the point, true and empty at once. Turning to a poem, they say it no more plainly, and yet the poem holds what the statement let slip. The poem did not close the gap the statement failed to close; it circled it, and in the circling let the unsayable thing be present as what the words turn around. What the letter could not deliver by aiming at it, the poem carries by not aiming at it.
\end{casebox}

\subsection*{The highest office of a language at its limit}

So the poem is the paradigm of a language reconciled to its limit and made generative by it. It does not mourn what it cannot say and does not pretend to say it. It builds its whole art on the not-saying, and out of the not-saying makes the one form of speech that keeps the gap open as a living center rather than closing it as a solved problem. What the poem shows, and shows in the medium of language itself, is that the limit of language is not the end of what language can do but the opening of its highest office. This bears directly on the two registers that follow. If the poem is the making of a circulation without capture in language, ethics and justice will be shown to be the making of a relation without possession, and a guarding without settlement, in the registers of the other and of the common life.
